\section{Supports and invariance under reduction}\label{sec:berkovich-reduction}

Let $k$ be a trivially valued field and let $A$ be a commutative $k$-algebra.
No finite generation, reducedness, or algebraic closedness assumption is needed in this section.
These affine foundations make explicit the support convention in \cite[\S2.1]{FHX} and the
convention in \S5 that closed subsets carry their reduced induced structure.

\begin{definition}[Support map {\cite[\S2.1]{FHX}}]\label{berk:support}
\lean{EnvelopeProperty.Berk.support, EnvelopeProperty.Berk.preimage_support_zeroLocus,
EnvelopeProperty.Berk.continuous_support, EnvelopeProperty.Berk.support_comap}
\leanok
\uses{def:berk, def:analytic-zero-locus, lem:analytic-zero-locus-closed,
spectrum:support, spectrum:trivial-field-bridge}
The support of $x\in(\Spec A)^\an$ is the prime ideal
$\mathfrak p_x=\{a\in A:|a|_x=0\}$. The resulting map
$\operatorname{supp}\colon(\Spec A)^\an\to\Spec A$ is continuous and satisfies
$\operatorname{supp}(\phi^\an(y))=\phi^{-1}(\operatorname{supp}(y))$ for every
$k$-algebra homomorphism $\phi\colon A\to B$. For any ideal $I$ of $A$,
$\operatorname{supp}^{-1}(V(I))=V(I)^\an$.
\end{definition}
\begin{proof}
The seminorm axioms imply that $\mathfrak p_x$ is an ideal and that
$ab\in\mathfrak p_x$ implies $a\in\mathfrak p_x$ or $b\in\mathfrak p_x$;
it is proper because $|1|_x=1$. A point lies over $V(I)$ exactly when every
element of $I$ has seminorm zero, proving the asserted inverse-image formula.
Each such zero locus is an intersection of the closed sets $\{|a|=0\}$.
As every Zariski closed set has the form $V(I)$, the support map is continuous.
Finally $|a|_{\phi^\an(y)}=|\phi(a)|_y$ proves the functoriality formula.
\end{proof}

\begin{lemma}[Every prime occurs as a support {\cite[\S2.1]{FHX}}]\label{berk:support-surjective}
\lean{EnvelopeProperty.Berk.primePoint, EnvelopeProperty.Berk.support_primePoint,
EnvelopeProperty.Berk.support_surjective}
\leanok
\uses{berk:support, def:triv-scaling, lem:analytic-affine-quotient}
For each prime ideal $\mathfrak p$ of $A$, the trivial seminorm of the domain
$A/\mathfrak p$ pulls back to an analytic point $x_{\mathfrak p}$ whose support
is $\mathfrak p$. Thus the support map is surjective.
\end{lemma}
\begin{proof}
The quotient is a domain, so its trivial seminorm is multiplicative and ultrametric.
Its restriction to the nonzero constants is one because the field $k$ embeds in
$A/\mathfrak p$. The pullback has $|a|_{x_{\mathfrak p}}=0$ exactly when
the image of $a$ in $A/\mathfrak p$ is zero, which is equivalent to $a\in\mathfrak p$.
\end{proof}

\begin{lemma}[Analytic vanishing and radicals]\label{berk:zero-locus-radical}
\lean{EnvelopeProperty.Berk.zeroLocus_radical,
EnvelopeProperty.Berk.zeroLocus_subset_zeroLocus_iff,
EnvelopeProperty.Berk.zeroLocus_eq_zeroLocus_iff}
\leanok
\uses{berk:support, berk:support-surjective}
For ideals $I,J$ of $A$, the analytic zero locus of $I$ equals that of its radical,
and $V(I)^\an\subseteq V(J)^\an$ if and only if $J\subseteq\sqrt I$.
Consequently $V(I)^\an=V(J)^\an$ if and only if $\sqrt I=\sqrt J$.
\end{lemma}
\begin{proof}
Because $\mathfrak p_x$ is prime, $I\subseteq\mathfrak p_x$ holds exactly when
$\sqrt I\subseteq\mathfrak p_x$, proving the first assertion.
The inverse-image formula for support and its surjectivity show that inclusion
of analytic zero loci is equivalent to inclusion of the corresponding Zariski zero loci.
The latter inclusion holds exactly when every prime containing $I$ contains $J$.
The intersection of all primes containing $I$ is $\sqrt I$, giving the inclusion criterion.
Applying that criterion in both directions gives equality of the two radicals.
\end{proof}

\begin{lemma}[Nilpotents are invisible to analytic points]\label{berk:nilpotents}
\lean{EnvelopeProperty.Berk.nilradical_le_supp,
EnvelopeProperty.Berk.val_eq_zero_of_isNilpotent,
EnvelopeProperty.Berk.val_eq_of_sub_isNilpotent,
EnvelopeProperty.Berk.logAbs_eq_of_sub_isNilpotent,
EnvelopeProperty.Berk.val_eq_zero_forall_iff_isNilpotent,
EnvelopeProperty.Berk.zeroLocus_nilradical}
\leanok
\uses{berk:support, berk:support-surjective, spectrum:radical-vanishing,
spectrum:reduction-homeomorphism}
The nilradical of $A$ is contained in the support of every analytic point.
In particular every nilpotent has seminorm zero at every point, and conversely
an element that has seminorm zero at every analytic point is nilpotent.
Functions differing by a nilpotent have equal seminorms and equal logarithmic evaluations.
The analytic zero locus of the nilradical is the whole analytification.
\end{lemma}
\begin{proof}
Every prime contains the nilradical, so every support contains it.
If $a$ has seminorm zero at every point, surjectivity of support implies that
$a$ belongs to every prime ideal. It therefore lies in their intersection,
the nilradical, whose elements are precisely the nilpotents.
If $a-b$ is nilpotent, then $|a-b|_x=0$. The ultrametric inequality applied to
$a=b+(a-b)$ gives $|a|_x\le|b|_x$, and applied to $b=a-(a-b)$ gives the reverse
inequality. Thus their seminorms and logarithms agree.
The final assertion follows immediately from the containment already proved.
\end{proof}

\begin{theorem}[Invariance under reduction {\cite[\S5]{FHX}}]\label{berk:reduction-homeomorphism}
\lean{EnvelopeProperty.Berk.reductionHomeomorph,
EnvelopeProperty.Berk.reductionHomeomorph_val_apply,
EnvelopeProperty.Berk.reductionHomeomorph_symm_val_mk}
\leanok
\uses{berk:nilpotents, lem:analytic-affine-quotient}
For $A_{\mathrm{red}}=A/\sqrt{(0)}$, restriction induces a homeomorphism
$(\Spec A_{\mathrm{red}})^\an\simeq(\Spec A)^\an$.
This homeomorphism preserves the seminorm of every $a\in A$ and the inverse
preserves the seminorm of its quotient class.
\end{theorem}
\begin{proof}
The quotient homeomorphism identifies $(\Spec A_{\mathrm{red}})^\an$ with the
subspace of points whose supports contain $\sqrt{(0)}$. By the preceding lemma,
this subspace is the whole analytification. The evaluation formulas for the
quotient homeomorphism give the asserted compatibility of seminorms.
\end{proof}
